\section{Primos auxiliares y control uniforme}
\label{sec:auxiliares}

La construcción necesita una pareja de primos que satisfaga a la vez
condiciones de descomposición y estimaciones analíticas en un intervalo
entero. El segundo primo se elige en una progresión determinada por el
primero. Esta nota desarrolla los mecanismos finitos y analíticos que
permiten hacer esas elecciones con un umbral común, y distingue las
estimaciones profundas que las alimentan.

\fuente{A, §3, pp.~7--12, proposiciones~3.1--3.2;
A93, §5, pp.~21--28, Proposición~5.1, Lema~5.2 y Proposición~5.3.
A conserva el ejemplar de 55 páginas; A93 es el de 93 páginas.
Las pruebas de segundo momento, máximo, selección e integración se
reconstruyen aquí bajo las entradas uniformes que se enuncian.
El cotejo Lean se delimita en §\ref{subsec:aux-lean}.}

\subsection{Dos estimaciones y un orden de elección}

Fijamos un entero $M\ge1$, un exponente $J\ge1$ y
$0<\varepsilon<1/52$. Para una escala real $Y$ grande ponemos
\begin{equation}
 X=Y^J,\qquad z=Y^{\varepsilon/5},\qquad
 z'=z(\log Y)^2,\qquad Z=X^{400}.
 \label{eq:aux-escalas}
\end{equation}
$X$ es el extremo del intervalo que debemos controlar y $Z$ es el
corte de una suma finita. Los primos se buscarán en intervalos de
extremos exponenciales en $z$ y $z'$. Todos los umbrales pueden
depender de $M,J,\varepsilon$, fijados antes de variar $Y$.

Para un primo $p\equiv3\pmod4$, escribimos $\chi_{-p}$ para el
carácter cuadrático primitivo de discriminante $-p$. Como función en
los enteros, $\chi_{-p}(n)=(n/p)$. La primera meta es elegir $q,d$
para que, para todo carácter $\xi$ módulo $M$ y toda pareja
$(i,j)\in\{0,1\}^2\setminus\{(0,0)\}$, el carácter primitivo
$\vartheta$ que induce $\xi\chi_{-q}^i\chi_{-d}^j$ satisfaga
\begin{equation}
 |G_\vartheta(y)|\le y^{1/2+\varepsilon/2},\qquad
 |\ell_\vartheta|\le4\log y
 \quad (Y\le y\le X).
 \label{eq:aux-meta-pareja}
\end{equation}
Los torcidos no triviales por los primos nuevos son no principales.
La cota para $G_\vartheta$ controla los caracteres auxiliares que se
retiran en la extracción de la Nota~\ref{sec:caracteres}.

La segunda meta se aplica a todos los caracteres del cuerpo anfitrión,
incluido el principal, pero en argumentos mucho mayores:
\begin{equation}
 y\ge\tfrac13\exp(x^{\varepsilon/3})
 \quad\Longrightarrow\quad
 |G_\vartheta(y)|\le
 y\exp\bigl(-\tfrac12x^{\varepsilon/8}\bigr).
 \label{eq:aux-meta-delgada}
\end{equation}
Aquí $x$ es la escala de la construcción y $y$ el argumento de la suma.
Una cota relativa en esta escala no es la cota de raíz cuadrada de
\eqref{eq:aux-meta-pareja} en $[Y,Y^J]$.

La selección que estudiaremos tiene el orden
\begin{equation}
 (M,J,\varepsilon)\ \longrightarrow\ Y_0\
 \longrightarrow\ Y\ge Y_0\ \longrightarrow\ q\
 \longrightarrow\ d\ \longrightarrow\
 \text{todos los caracteres y todos los }y\in[Y,Y^J].
 \label{eq:aux-orden}
\end{equation}
El umbral de la segunda etapa debe servir para todos los primeros
primos admisibles. Un umbral elegido después de cada $q$ no daría
la uniformidad necesaria al ensamblaje.

\subsection{Congruencias, descomposición y módulo creciente}

Definimos conjuntos finitos de primos
\begin{align}
 \mathcal P_1
 &=\{q:e^{z/2}<q\le e^z,\ q\equiv3\pmod8\},\notag\\
 R_q&=q\prod_{\substack{3\le p\le z\\p\text{ primo}}}p,\notag\\
 \mathcal P_2(q)
 &=\{d:e^{z'/2}<d\le e^{z'},\ d\equiv7\pmod{16},\
                                      R_q\mid d+1\}.
 \label{eq:aux-conjuntos}
\end{align}
Para $Y$ grande, $q>z$ y $R_q$ es impar y libre de cuadrados.
También podemos exigir $q>M$ y $d>q$; así los dos primos son distintos
y coprimos a $M$.

\begin{lema}[La progresión del segundo primo]
\label{lem:aux-crt}
Si $R$ es impar, las condiciones $d\equiv7\pmod{16}$ y
$d\equiv-1\pmod R$ equivalen a
\begin{equation}
 d\equiv8R-1\pmod{16R}.
 \label{eq:aux-crt}
\end{equation}
Esta clase es coprima a $16R$.
\end{lema}

\begin{proof}
Como $R$ es impar, $8R-1\equiv7\pmod{16}$ y
$8R-1\equiv-1\pmod R$. El teorema chino del resto da una única
clase módulo $16R$, pues $\gcd(16,R)=1$. Su representante es impar
y no comparte un primo con $R$ porque es $-1$ módulo $R$.
\end{proof}

Las congruencias tienen un uso aritmético concreto. El criterio clásico
para un discriminante cuadrático impar dice que $2$ es inerte si el
discriminante es $5$ módulo $8$, y se descompone si es $1$ módulo $8$.
Por ello $2$ es inerte en $\Q(\sqrt{-q})$ y se descompone en
$\Q(\sqrt{-d})$. Para cada primo impar $p\le z$, y también para
$p=q$, se tiene $-d\equiv1\pmod p$, luego $(-d/p)=1$ y esos primos
se descomponen en $\Q(\sqrt{-d})$. Estos son los criterios usados
por A93, p.~27; la construcción geométrica posterior recibe tales datos.

La cota clásica de Chebyshev para
$\theta(z)=\sum_{p\le z}\log p$ permite tomar $\theta(z)\le2z$
para $z$ suficientemente grande. Así, uniformemente en $q\in\mathcal P_1$,
\begin{equation}
 \log(16R_q)\le\log16+z+\theta(z)\le4z,
 \qquad \tfrac14(\log Y)^2\log(16R_q)\le z'.
 \label{eq:aux-modulo}
\end{equation}
El módulo de la progresión crece con $Y$. De hecho,
$\log(16R_q)>z/2$, mientras $\sqrt{z'}=\sqrt z\log Y$, de modo
que $\log(16R_q)/\sqrt{z'}$ tiende a infinito. Un resultado que
solo controle módulos fijos, o que exija un módulo mucho menor en
esta comparación, no proporciona el conteo que se necesita.

\subsection{Las entradas analíticas que se reciben}
\label{subsec:aux-entradas}

Sea
\[
 \mathcal P=\{p:\sqrt U<p\le U,\ p\equiv a_0\pmod{m_0}\},
 \qquad \gcd(a_0,m_0)=1,\quad8\mid m_0.
\]
El rango de progresiones de A93, (5.1), es
\begin{equation}
 \tfrac14(\log Y)^2\log m_0\le\log U,
 \qquad Y^{\varepsilon/6}\le\log U
                  \le Y^{\varepsilon/5}(\log Y)^2.
 \label{eq:aux-rango}
\end{equation}
Se reciben las siguientes estimaciones, con uniformidad en los datos
variables de ese rango.

\begin{enumerate}[label=\textbf{\arabic*.}]
\item \textbf{Conteo de primos.} Para $Y$ suficientemente grande,
$|\mathcal P|\ge U^{9/10}$.
\item \textbf{Cancelación salvo una clase cuadrática.} Dados exponentes
fijos $A,B>0$, existe a lo sumo un entero impar libre de cuadrados
$s_*>1$, coprimo a $m_0$, que hay que excluir para tener
\begin{equation}
 \left|\sum_{p\in\mathcal P}\left(\frac p s\right)\right|
 \le Y^{-B}|\mathcal P|
 \quad\left(1<s\le Y^A,\ s\text{ impar libre de cuadrados},\
 \gcd(s,m_0)=1,\ s\ne s_*\right).
 \label{eq:aux-cancelacion}
\end{equation}
\item \textbf{Densidad de ceros.} El número de caracteres primitivos
de conductor a lo sumo $Q$ con un cero en
$\Ree s>399/400$, $|\Im s|\le Q$, es $O(Q^{1/4})$.
\item \textbf{Cola y derivada logarítmica.} Una caja sin ceros con
los rangos de conductor indicados más adelante proporciona
\eqref{eq:aux-cola} y \eqref{eq:aux-loglog}.
\item \textbf{Fórmula explícita uniforme.} Tras controlar los ceros
reales de los caracteres pertinentes, se dispone de
\eqref{eq:aux-psi}, con el término polar principal sustraído.
\end{enumerate}

Las tres primeras entradas corresponden a \texttt{MSqCount},
\texttt{twistK} y \texttt{CountBad}; las otras se precisan al usarlas.
A93, pp.~22--23, conserva el posible cero excepcional dentro del término
principal del conteo de progresiones y estima el error relativamente
a ese término. El teorema de Siegel interviene con una constante
posiblemente inefectiva. Esto afecta al cálculo del umbral, no al orden
lógico de la existencia. Esta nota no reconstruye esos teoremas uniformes
a partir de ortogonalidad elemental.

\subsection{La clase excepcional en el segundo momento}

Para $y\ge3$, $Z\ge y$ y una secuencia compleja común
$|\kappa(n)|\le1$, ponemos
\begin{equation}
 c_y(n)=\Lambda(n)\min(1,y/n),\qquad
 S_p(y)=\sum_{n\le Z}c_y(n)\kappa(n)\left(\frac n p\right).
 \label{eq:aux-suma-truncada}
\end{equation}
$S_p$ designa aquí una suma truncada sobre el primo que se está eligiendo;
no es la suma suavizada $S_\chi$ de la Nota~\ref{sec:sumas}.
La secuencia $\kappa$ puede cambiar entre aplicaciones, pero no depende
del primo $p$ sobre el que se promedia en una misma aplicación.

Suponemos $p>Z^2$ para todo $p\in\mathcal P$.
Para cada entero $n\ge1$, sea $s(n)$ el producto de los primos que no dividen
$m_0$ y aparecen con exponente impar en $n$.
Como $c_y$ está soportada en potencias de primos, $s(n)$ vale $1$
o un solo primo. Definimos
\begin{equation}
 R=\sum_{\substack{n\le Z\\s(n)=1}}c_y(n),\qquad
 A_\ell=\sum_{\substack{k\ge1\text{ impar}\\\ell^k\le Z}}
                         c_y(\ell^k)\ (\ell\nmid m_0),\qquad
 B_0=\sup_{\ell\nmid m_0}A_\ell.
 \label{eq:aux-coeficientes}
\end{equation}
Solo un número finito de $A_\ell$ es distinto de cero; si todos son cero,
tomamos $B_0=0$. La letra $R$ en esta subsección es una suma de pesos,
distinta del entero $R_q$ de \eqref{eq:aux-conjuntos}.

\begin{lema}[Peso de las clases cuadráticas]
\label{lem:aux-clases}
La suma de los pesos $c_y(n)c_y(n')$ de los pares cuya parte cuadrática
exterior es $1$ vale $R^2+\sum_\ell A_\ell^2$.
Para una clase fija $s>1$, esa suma es a lo sumo
$2B_0R+2B_0^2$. Bajo \eqref{eq:aux-cancelacion}, si $Z^2\le Y^A$,
\begin{equation}
 \frac1{|\mathcal P|}\sum_{p\in\mathcal P}|S_p(y)|^2
 \le R^2+\sum_\ell A_\ell^2+2B_0R+2B_0^2
       +Y^{-B}\left(\sum_{n\le Z}c_y(n)\right)^2.
 \label{eq:aux-momento-finito}
\end{equation}
Se supone $\mathcal P\ne\varnothing$ en esta fórmula.
\end{lema}

\begin{proof}
La parte cuadrática exterior de $nn'$ se obtiene multiplicando
$s(n)$ y $s(n')$ y retirando los factores cuadrados. Es $1$ exactamente
cuando ambos valores son $1$ o son el mismo primo, lo que da
$R^2+\sum_\ell A_\ell^2$.
Si $s=\ell$ es primo, los pares posibles tienen partes exteriores
$1,\ell$ en uno de los dos órdenes y pesan $2RA_\ell$.
Si $s=\ell r$ con $\ell\ne r$ primos, pesan $2A_\ell A_r$.
Si $s$ tiene tres o más factores primos, no hay pares. Esto da la
cota uniforme anunciada para cualquier clase excepcional.

Al expandir $|S_p|^2$ aparece
$\kappa(n)\overline{\kappa(n')}$, de módulo a lo sumo uno.
La reciprocidad cuadrática y la clase fija de $p$ módulo $m_0$ dan
\[
 \left(\frac n p\right)\left(\frac{n'}p\right)
 =\epsilon_{n,n'}\left(\frac p{s(nn')}\right),
 \qquad \epsilon_{n,n'}\in\{-1,1\},
\]
con signo independiente de $p$. La condición $8\mid m_0$ fija
las contribuciones de $2$ y de los signos de reciprocidad; los primos
que dividen $m_0$ aportan también signos fijos. $p>Z^2$ evita símbolos
nulos. Se estima trivialmente la clase $1$ y la posible clase $s_*$;
para las demás se usa \eqref{eq:aux-cancelacion}. Sus pesos totales
son a lo sumo $(\sum c_y)^2$, lo que prueba la desigualdad.
\end{proof}

Los términos mixtos $2B_0R$ importan: una potencia impar de un primo
exterior y una potencia par de otro primo pueden pertenecer a una clase
excepcional prima. El argumento conserva la clase excepcional y controla
su peso, sin necesitar conocer cuál es.

\subsection{De los coeficientes a una cota uniforme}

\begin{proposicion}[Segundo momento de la suma truncada]
\label{prop:aux-momento}
Fijados los datos de \eqref{eq:aux-escalas}, supongamos que todo primo
que divide $m_0$ es a lo sumo $z$ o es mayor que $Z$, y que se cumple
\eqref{eq:aux-cancelacion} con $A=800J$ y $B=2J+2$.
Entonces existe una constante $C_{\mathrm{ms}}$, uniforme en la secuencia
$|\kappa|\le1$ y en los demás datos admitidos, tal que
\begin{equation}
 \sum_{p\in\mathcal P}|S_p(y)|^2
 \le C_{\mathrm{ms}}|\mathcal P|\,y\log(2y)
 \qquad(Y\le y\le X).
 \label{eq:aux-momento}
\end{equation}
Se conservan las hipótesis de primalidad y tamaño del lema anterior.
\end{proposicion}

\begin{proof}
Si $\mathcal P$ es vacío la cota es inmediata; supongámoslo no vacío.
Usamos la cota clásica de Chebyshev
$\Psi(u)=\sum_{n\le u}\Lambda(n)\ll u$.
Para un primo $\ell\le y$, sus potencias hasta $y$ aportan a lo sumo
$\log y$; la cola geométrica de las potencias mayores aporta
$O(\log\ell)$. Para $\ell>y$, todas aportan
$O(y\log\ell/\ell)$. Se deduce $B_0\ll\log(2y)$.

Las potencias pares satisfacen
\begin{equation}
 \sum_{\ell,k\ge1}\log\ell\min(1,y/\ell^{2k})
 =2y\int_{\sqrt y}^{\infty}\frac{\Psi(t)}{t^3}\dd t
 \ll\sqrt y.
 \label{eq:aux-pares}
\end{equation}
La identidad se obtiene escribiendo cada mínimo como la integral del
núcleo $2y/t^3$ desde $\max(\sqrt y,\ell^k)$ hasta infinito;
la no negatividad permite sumar. No se deriva la función escalonada
$\Psi$. Los primos que dividen $m_0$ y son mayores que $Z$ no
contribuyen; los restantes añaden $O(z\log(2y))$. Como
$z\log(2y)/\sqrt y\le Y^{\varepsilon/5-1/2}
(\log2+J\log Y)$, se obtiene $R\ll\sqrt y$ uniformemente.

Para $\ell\le y$, la suma de los $A_\ell$ es $O(y)$: las
potencias hasta $y$ se acotan por $\Psi(y)$ y las colas por
$\sum_{\ell\le y}\log\ell\ll y$.
Por ello $\sum_{\ell\le y}A_\ell^2\ll y\log(2y)$.
Para los primos mayores que $y$, sumación parcial da
\[
 \sum_{\ell>y}A_\ell^2
 \ll y^2\sum_{\ell>y}\frac{(\log\ell)^2}{\ell^2}
 \ll y\log(2y).
\]
En la última estimación se aplica la cota de Chebyshev al factor
$\log\ell$ y al núcleo decreciente $(\log t)/t^2$.

También por sumación parcial,
\[
 \sum_{n\le Z}c_y(n)
 =\frac{y\Psi(Z)}Z+y\int_y^Z\frac{\Psi(t)}{t^2}\dd t
 \ll y(1+\log Z).
\]
Como $Z=Y^{400J}$, la última contribución de
\eqref{eq:aux-momento-finito} está acotada por
\[
 Y^{-2J-2}X^2(1+400J\log Y)^2
 =Y^{-2}(1+400J\log Y)^2=O_J(1).
\]
Además $Z^2=Y^{800J}$, como exige el rango de cancelación.
Sustituyendo estas cotas en \eqref{eq:aux-momento-finito} se obtiene
\eqref{eq:aux-momento}. Los términos $B_0R$ y $B_0^2$ también
quedan absorbidos por $y\log(2y)$.
\end{proof}

La uniformidad en $\kappa$ es más fuerte que una afirmación solo
para caracteres multiplicativos. Será necesaria para controlar la
derivada de la suma truncada.

\subsection{Un lema maximal para todo el intervalo}

\begin{lema}[Control simultáneo a partir de dos momentos]
\label{lem:aux-maximal}
Sea $\mathcal P$ un conjunto finito y $2\le Y\le X$.
Para cada $p\in\mathcal P$, sea $F_p:[Y,X]\to\C$
absolutamente continua. Supongamos que, para una constante $C\ge0$,
\begin{align}
 \sum_p|F_p(t)|^2&\le C|\mathcal P|t\log(2t)
                         &&\text{para todo }t,\notag\\
 \sum_p|tF'_p(t)|^2&\le C|\mathcal P|t\log(2t)
                         &&\text{para casi todo }t.
 \label{eq:aux-dos-momentos}
\end{align}
Dados $h,\tau>0$, el número de índices para los que existe
$t\in[Y,X]$ con $|F_p(t)|>ht^{1/2+\tau}$ es a lo sumo
\begin{equation}
 C h^{-2}|\mathcal P|Y^{-2\tau}
 \left[\left(1+\frac1\tau\right)\log(2Y)
                         +\frac1{2\tau^2}\right].
 \label{eq:aux-maximal}
\end{equation}
La cota no contiene $X$.
\end{lema}

\begin{proof}
Si $\mathcal P$ es vacío la afirmación es inmediata.
Pongamos $Q_p(t)=t^{-1-2\tau}|F_p(t)|^2$.
Esta función es absolutamente continua y, casi en todas partes,
\begin{align*}
 Q'_p(t)
 &=2t^{-1-2\tau}\Ree\bigl(F'_p(t)\overline{F_p(t)}\bigr)
       -(1+2\tau)t^{-2-2\tau}|F_p(t)|^2\\
 &\le2t^{-2-2\tau}|F_p(t)|\,|tF'_p(t)|.
\end{align*}
El segundo término es no positivo. Integrando hasta cualquier
$v\in[Y,X]$ y tomando el supremo,
\[
 \sup_{Y\le v\le X}Q_p(v)
 \le Q_p(Y)+2\int_Y^X
             \frac{|F_p(t)|\,|tF'_p(t)|}{t^{2+2\tau}}\dd t.
\]
Sumamos en $p$ y aplicamos Cauchy--Schwarz a las dos familias de
\eqref{eq:aux-dos-momentos}. Resulta
\begin{equation}
 \sum_p\sup_{Y\le v\le X}Q_p(v)
 \le C|\mathcal P|Y^{-2\tau}\log(2Y)
       +2C|\mathcal P|\int_Y^\infty
                    t^{-1-2\tau}\log(2t)\dd t.
 \label{eq:aux-supremos}
\end{equation}
Solo se ha prolongado hasta infinito una integral comparadora positiva;
las funciones $F_p$ siguen definidas únicamente en $[Y,X]$.
Por integración por partes,
\begin{equation}
 \int_Y^\infty t^{-1-2\tau}\log(2t)\dd t
 =Y^{-2\tau}\left(\frac{\log(2Y)}{2\tau}
                                     +\frac1{4\tau^2}\right).
 \label{eq:aux-integral-maximal}
\end{equation}
Cada índice excepcional aporta más de $h^2$ a la suma de supremos.
Dividir \eqref{eq:aux-supremos} por $h^2$ y sustituir la integral
prueba \eqref{eq:aux-maximal}.
\end{proof}

La continuidad absoluta permite recuperar cada valor de $Q_p$ a partir
de su derivada integrable. Una cota para $F_p$ sin control de la
variación no impediría que distintos índices tuvieran picos en
puntos diferentes, hasta cubrir todos los índices al recorrer el intervalo.

\subsection{La derivada y la familia de caracteres}

Cada $S_p$ de \eqref{eq:aux-suma-truncada} es una suma finita de
funciones continuas y afines a trozos, luego es absolutamente continua.
Fuera de sus puntos de quiebre,
\begin{equation}
 tS'_p(t)=\sum_{t<n\le Z}\Lambda(n)\kappa(n)
                         \left(\frac n p\right)\frac t n.
 \label{eq:aux-derivada}
\end{equation}
Es la misma suma de pesos $c_t(n)$ con la secuencia
$\kappa_t(n)=\kappa(n)\mathbf1_{n>t}$. Esta secuencia tiene módulo
a lo sumo uno y es común a todos los primos promediados. Aunque depende
de $t$ y puede perder multiplicatividad, la Proposición~\ref{prop:aux-momento}
se aplica con la misma constante. Se obtienen así los dos momentos
necesarios para el lema maximal.

\begin{proposicion}[Una familia finita con el mismo umbral]
\label{prop:aux-familia}
Sea $\mathcal P\ne\varnothing$. Supongamos que para $r$ familias de funciones se cumplen
\eqref{eq:aux-dos-momentos} con constantes $C_1,\ldots,C_r$.
La proporción de índices que falla en alguna familia o algún punto
está acotada por
\begin{equation}
 h^{-2}\left(\sum_{a=1}^r C_a\right)D_\tau(Y),
 \qquad
 D_\tau(Y)=Y^{-2\tau}
 \left[\left(1+\frac1\tau\right)\log(2Y)+\frac1{2\tau^2}\right].
 \label{eq:aux-union}
\end{equation}
En particular, si $r\le r_0$ y $C_a\le C_0$ uniformemente, para
cada $\eta>0$ hay un umbral común que hace esa proporción menor
que $\eta$.
\end{proposicion}

\begin{proof}
La cardinalidad de una unión finita es a lo sumo la suma de las
cardinalidades. Aplicamos \eqref{eq:aux-maximal} a cada familia y
dividimos por $|\mathcal P|$, suponiéndolo positivo.
Como $\tau>0$ es fijo, $D_\tau(Y)\to0$; basta imponer
$h^{-2}r_0C_0D_\tau(Y)<\eta$.
\end{proof}

En la primera etapa se usan los caracteres primitivos que inducen
caracteres módulo $8M$. En la segunda, para un $q$ fijado, se
permiten además sus productos con $\chi_{-q}$: hay a lo sumo
$2\varphi(8M)$ funciones base. Sus conductores varían, pero su
módulo sigue acotado por uno, que es la propiedad utilizada en
\eqref{eq:aux-momento}. El número de condiciones no crece con $q$.
Se obtiene un umbral común a todas las fibras; no se hace una unión
sobre todos los primeros primos posibles.

La independencia de $X$ en \eqref{eq:aux-maximal} presupone un
$C$ fijo. En su aplicación, la cancelación y la truncación usan
$A=800J$, $B=2J+2$ y $Z=Y^{400J}$. El umbral completo puede
por tanto depender de $J$, como permite \eqref{eq:aux-orden}.

\subsection{Conductores primitivos y exclusión por densidad de ceros}

Un carácter módulo $M$ puede ser imprimitivo. Su carácter primitivo
inductor tiene los mismos valores en los enteros coprimos a $M$, pero
puede tener valores distintos en los primos que dividen $M$ y no su
conductor. En esta nota se toma primero el inductor primitivo y se
utiliza esa función como $\kappa$. Esta elección permite aplicar a la
misma suma la identidad con $G_\vartheta$; no se suprimen factores de
Euler sin registrarlos.

\begin{lema}[El conductor identifica el primo nuevo]
\label{lem:aux-inyeccion}
Sean $1\le F\le\sqrt U$ y $\mathcal P$ un conjunto de primos
$p\equiv3\pmod4$ contenido en $(\sqrt U,U]$.
Para cada $p\in\mathcal P$,
si $\kappa$ es primitivo de conductor $f\le F$, entonces
$\kappa\chi_{-p}$ es primitivo de conductor $fp$.
Si se reciben el conteo $|\mathcal P|\ge U^{9/10}$ y la entrada de
densidad de §\ref{subsec:aux-entradas}, la proporción de primos para
los que algún carácter de esta forma tiene un cero en la caja de
altura $FU$ está acotada por
\begin{equation}
 O(U^{-21/40}).
 \label{eq:aux-densidad}
\end{equation}
\end{lema}

\begin{proof}
Los conductores $f$ y $p$ son coprimos. Por el teorema chino del resto,
el carácter producto sobre las unidades módulo $fp$ se restringe
a $\kappa$ y a $\chi_{-p}$ en los dos factores. Si su conductor
fuera un divisor propio de $fp$, una de esas restricciones se
induciría desde un divisor propio del conductor correspondiente,
contradiciendo su primitividad. Por tanto el conductor es $fp$.

Para cada primo malo elegimos un carácter que exhiba el fallo.
En su conductor, $p$ es el único factor primo mayor que $F$,
pues todos los factores de $f$ son a lo sumo $F$. Así un mismo
carácter no puede elegirse para dos primos nuevos distintos,
incluso si el carácter base cambia. Esta aplicación es inyectiva
en el conjunto de caracteres excepcionales de conductor a lo sumo
$FU$ y caja de altura $FU$. Su cardinalidad es
$O((FU)^{1/4})\le O(U^{3/8})$.
Dividir por $U^{9/10}$ da el exponente
$3/8-9/10=-21/40$. No hay que multiplicar por el número de posibles
caracteres base en esta cuenta inyectiva.
\end{proof}

\subsection{Elegir el segundo primo en su propia fibra}

\begin{lema}[Selección sucesiva]
\label{lem:aux-fibras}
Sea $\mathcal Q$ un conjunto finito no vacío y, para cada $q\in\mathcal Q$,
sea $\mathcal D_q$ un conjunto finito no vacío. Fijemos $0<\eta<1$.
Supongamos que $\mathcal Q_0\subseteq\mathcal Q$ y
$\mathcal D_q^0\subseteq\mathcal D_q$ para $q\in\mathcal Q_0$ cumplen
\begin{equation}
 |\mathcal Q_0|\ge(1-\eta)|\mathcal Q|,\qquad
 |\mathcal D_q^0|\ge(1-\eta)|\mathcal D_q|.
 \label{eq:aux-fibras}
\end{equation}
Existe una pareja buena $(q,d)$ con $q\in\mathcal Q_0$ y
$d\in\mathcal D_q^0$. Si se elige $q$ uniformemente en $\mathcal Q$
y después $d$ uniformemente en su fibra $\mathcal D_q$, la probabilidad
de obtener una pareja buena es al menos $(1-\eta)^2$.
\end{lema}

\begin{proof}
Las dos cotas inferiores son estrictamente positivas; sus cardinalidades
enteras dan sucesivamente un $q$ bueno y un $d$ bueno en su fibra.
Para la afirmación probabilística,
\[
 \mathbb P(\text{pareja buena})
 =\frac1{|\mathcal Q|}\sum_{q\in\mathcal Q_0}
             \frac{|\mathcal D_q^0|}{|\mathcal D_q|}
 \ge\frac{|\mathcal Q_0|}{|\mathcal Q|}(1-\eta)
 \ge(1-\eta)^2.
\]
No se ha supuesto independencia ni igualdad de tamaños de las fibras.
\end{proof}

En la aplicación, $\mathcal D_q=\mathcal P_2(q)$.
El techo común de conductores de la segunda etapa es
\begin{equation}
 F=8Me^z,\qquad \log F=\log(8M)+z=o(z'),\qquad
 F<e^{z'/2}=\sqrt U
 \label{eq:aux-techo}
\end{equation}
para $Y$ suficientemente grande. Sirve para todos los $q\le e^z$.
La cota de módulo \eqref{eq:aux-modulo} y la separación
$e^{z/2}>Z^2=Y^{800J}$ también se obtienen con umbrales fijados
antes de elegir $q$. El máximo de esos umbrales permite aplicar
la selección sucesiva en una misma escala.

La probabilidad del lema corresponde al muestreo indicado. Si las fibras
tienen tamaños diferentes, la medida uniforme sobre todos los pares
$(q,d)$ asigna otros pesos a los primeros primos; no debe confundirse
con este muestreo sucesivo.

\subsection{El signo de la suma truncada y la cola}

\begin{lema}[Identidad exacta de truncación]
\label{lem:aux-truncacion}
Sea $\vartheta$ primitivo no principal y $Z\ge y\ge1$. Definamos
\begin{equation}
 T_\vartheta(Z)=-\ell_\vartheta-
                  \sum_{n\le Z}\frac{\Lambda(n)\vartheta(n)}n.
 \label{eq:aux-t}
\end{equation}
Entonces
\begin{equation}
 G_\vartheta(y)
 =-\sum_{n\le Z}\Lambda(n)\vartheta(n)\min(1,y/n)
                                      -yT_\vartheta(Z).
 \label{eq:aux-truncacion}
\end{equation}
\end{lema}

\begin{proof}
El primer término del lado derecho es
$-\sum_{n\le y}\Lambda(n)\vartheta(n)
-y\sum_{y<n\le Z}\Lambda(n)\vartheta(n)/n$.
Al sumar $-yT_\vartheta(Z)$, se cancelan los términos con $n>y$;
quedan $\sum_{n\le y}\Lambda(n)\vartheta(n)(y/n-1)
+y\ell_\vartheta$, que es la definición de $G_\vartheta$.
\end{proof}

La identidad tiene signo negativo delante de la suma truncada.
Para un torcido $\vartheta=\kappa\chi_{-p}$, esa suma es el
$S_p$ de \eqref{eq:aux-suma-truncada} con la base primitiva apropiada.
La entrada de contorno \texttt{TailKExp}, usada por \texttt{BoxOfS},
proporciona
\begin{equation}
 X|T_\vartheta(X^{400})|\le Y^\varepsilon
 \label{eq:aux-cola}
\end{equation}
cuando el conductor es a lo sumo $Q$, la caja especificada está libre
de ceros y $Y^{\varepsilon/6}\le\log Q\le Y^{\varepsilon/4}$.
Su reserva de exponente depende de $J$ mediante
$\varepsilon/(1600J)$; no se obtiene una conclusión uniforme en
$J$ eliminando esa dependencia.

Aplicamos el lema maximal con $h=1/2$ y $\tau=\varepsilon/2$.
Para los primos buenos, $|S_p(y)|\le\tfrac12 y^{1/2+\varepsilon/2}$
en todo el intervalo. Además,
\[
 y|T_\vartheta(X^{400})|\le Y^\varepsilon
       \le\tfrac12 y^{1/2+\varepsilon/2}
\]
para $Y$ grande, porque $1/2-\varepsilon/2>0$.
La desigualdad triangular en \eqref{eq:aux-truncacion} da la primera
cota de \eqref{eq:aux-meta-pareja}.

La entrada \texttt{Kappa4Loglog} controla también
\begin{equation}
 |\ell_\vartheta|\le1000\log\log Q+C.
 \label{eq:aux-loglog}
\end{equation}
Aquí $C$ es uniforme en el carácter seleccionado. Si
$\log Q\le C_{M,J}Y^{\varepsilon/5}(\log Y)^2$, resulta
\[
 |\ell_\vartheta|\le200\varepsilon\log Y
                    +2000\log\log Y+O_{M,J}(1).
\]
Como $200\varepsilon<4$, los términos de orden inferior se absorben
en la holgura y se obtiene $|\ell_\vartheta|\le4\log Y\le4\log y$.
La dependencia doblemente logarítmica del conductor es decisiva en
esta cuenta; una cota por un múltiplo de $\log Q$ sería demasiado grande.

\subsection{Ensamblaje condicionado de la pareja}

\begin{proposicion}[Pareja con estimaciones en todo el intervalo]
\label{prop:aux-pareja}
Bajo las entradas uniformes de §\ref{subsec:aux-entradas},
\eqref{eq:aux-cola} y \eqref{eq:aux-loglog}, existe $Y_0$, dependiente
de los datos fijados en \eqref{eq:aux-escalas}, tal que para todo
$Y\ge Y_0$ hay $q\in\mathcal P_1$ y $d\in\mathcal P_2(q)$,
coprimos a $M$, que satisfacen \eqref{eq:aux-meta-pareja} para todos
los caracteres y torcidos allí indicados. Para cada carácter real
primitivo $\rho$ de conductor divisor de $8M$, los torcidos no triviales
por $\chi_{-q}$ y $\chi_{-d}$ no tienen ceros en
\begin{equation}
 \Ree s>399/400,\qquad |\Im s|\le Q_j,
 \qquad Q_j=8M\exp(z+jz'),\quad j\in\{0,1\},
 \label{eq:aux-cajas}
\end{equation}
con $j$ igual al exponente del torcido por $\chi_{-d}$.
\end{proposicion}

\begin{proof}
En la primera etapa tomamos $U=e^z$, $m_0=8$ y la clase $3$.
En la segunda tomamos $U=e^{z'}$, $m_0=16R_q$ y la clase
\eqref{eq:aux-crt}. La cota \eqref{eq:aux-modulo} coloca ambas
etapas en \eqref{eq:aux-rango} al aumentar un único umbral.
Los conjuntos son no vacíos por el conteo recibido.

En ambos casos los primos promediados son mayores que $Z^2$.
Los factores primos del segundo módulo son a lo sumo $z$, salvo
$q>Z$; el primer módulo solo tiene el factor $2$.
La Proposición~\ref{prop:aux-momento} y el lema maximal controlan
las sumas truncadas de cada base. La Proposición~\ref{prop:aux-familia}
reúne las condiciones para las bases módulo $8M$, y sus a lo sumo
$2\varphi(8M)$ variantes en la segunda etapa, con coste uniforme.

El Lema~\ref{lem:aux-inyeccion} descarta una proporción que tiende
a cero para las cajas de ceros. En la primera etapa se usa el techo
$8M$ y en la segunda el techo $F=8Me^z$, válido uniformemente por
\eqref{eq:aux-techo}. Los techos finales de conductor y altura son
precisamente $Q_0$ y $Q_1$. Sus logaritmos satisfacen los rangos de
\eqref{eq:aux-cola}, pues
$\log Q_1=\log(8M)+z+z'$ y
$(\log Y)^2=o(Y^{\varepsilon/20})$.

Fijamos, por ejemplo, una pérdida total menor que $1/2$ en cada etapa.
Se distribuye esa pérdida entre los conjuntos excepcionales finitos y
se toma el máximo de sus umbrales antes de elegir los primos.
El Lema~\ref{lem:aux-fibras} da una pareja compatible.
La identidad de truncación y \eqref{eq:aux-cola} producen la cota
para $G$; \eqref{eq:aux-loglog} produce la de $\ell$.
Tomar además $q>M$ y $d>q$ asegura la coprimalidad y distinción.
Los caracteres módulo $M$ quedan incluidos al trabajar con las bases
módulo $8M$, y los reales de conductor divisor de $8M$ quedan
incluidos en la exclusión de cajas. Esto prueba todas las afirmaciones.
\end{proof}

La proposición reconstruye el ensamblaje de A93, Proposición~5.1,
a partir de entradas analíticas delimitadas. No constituye aquí una
nueva demostración de las estimaciones de progresiones, densidad de ceros
ni contorno que se han recibido.

\subsection{La escala delgada y la suma con el polo sustraído}

Fijemos $B,D>0$ y un entero $Q_0\ge1$.
Consideremos caracteres primitivos de conductor divisor
de $Q_0qd$, donde $q,d$ son primos impares distintos y coprimos a $Q_0$,
la parte impar de $Q_0$ es a lo sumo $B$ y
\begin{equation}
 \log(Q_0qd)\le D x^{\varepsilon/5}(\log x)^2.
 \label{eq:aux-conductor-delgado}
\end{equation}
Suponemos que los caracteres reales no principales cuyo conductor
contiene $q$ o $d$ no tienen ceros reales mayores que $399/400$.

Un carácter real primitivo no principal tiene conductor de discriminante fundamental:
su parte impar es libre de cuadrados y su parte de potencia de $2$
divide $8$. Los que no contienen $q$ ni $d$ tienen conductor a lo
sumo $8B$ y forman una familia finita. La no anulación de sus funciones
$L$ en $1$ da un intervalo real común libre de ceros a la izquierda de
$1$. Para los otros caracteres se utiliza la hipótesis anterior.
Esta reducción permite aplicar la región clásica sin ceros y la entrada
uniforme de fórmula explícita \texttt{PsiThin}.

Escribamos
\[
 \Psi_0(\vartheta,u)=\sum_{n\le u}\Lambda(n)\vartheta(n)
                              -\mathbf1_{\vartheta=1}u,
 \qquad \lambda=\max(1,\log\operatorname{cond}\vartheta).
\]
La entrada que se utiliza, A93, (5.8), es
\begin{equation}
 |\Psi_0(\vartheta,u)|
 \le C u(\log u+\lambda)^2
      \exp\left(-\frac{c\log u}{\lambda+\sqrt{\log u}}\right),
 \qquad u\ge e,
 \label{eq:aux-psi}
\end{equation}
con $c,C>0$ uniformes en la familia. El carácter principal se incluye
con su término polar sustraído. En lo que sigue se demuestra el paso
de esta estimación a \eqref{eq:aux-meta-delgada}.

\begin{lema}[Representación integral de la suma corregida]
\label{lem:aux-integral-g}
Bajo \eqref{eq:aux-psi}, para $y\ge1$ se tiene
\begin{equation}
 G_\vartheta(y)=-y\int_y^\infty
                         \frac{\Psi_0(\vartheta,u)}{u^2}\dd u.
 \label{eq:aux-integral-g}
\end{equation}
La fórmula vale también para el carácter primitivo principal, con la
normalización de \eqref{eq:g-principal}.
\end{lema}

\begin{proof}
Fijemos el carácter. Para $v=\log u$ suficientemente grande,
$\lambda\le\sqrt v$, y \eqref{eq:aux-psi} prueba la convergencia
absoluta de la integral mediante una cota por un múltiplo de
$\int^\infty v^2e^{-c\sqrt v/2}\dd v$.
Por sumación parcial, para $s>1$ real,
\[
 s\int_1^\infty\frac{\Psi_0(\vartheta,u)}{u^{s+1}}\dd u
 =-\frac{L'(s,\vartheta)}{L(s,\vartheta)}
                  -\mathbf1_{\vartheta=1}\frac{s}{s-1}.
\]
La convergencia absoluta recién probada permite pasar al límite
$s\downarrow1$ por convergencia dominada. Para un carácter no principal
el límite es $-\ell_\vartheta$.
Para el principal, la expansión
$-\zeta'(1+t)/\zeta(1+t)=1/t-\gamma+O(t)$ y
$(1+t)/t=1/t+1$ da el límite $-1-\gamma$.
Así
\[
 \int_1^\infty\frac{\Psi_0(\vartheta,u)}{u^2}\dd u
 =\begin{cases}-\ell_\vartheta,&\vartheta\ne1,\\
               -1-\gamma,&\vartheta=1.\end{cases}
\]
La suma suavizada de \eqref{eq:suma-suave} satisface
$S_\vartheta(y)=y\int_1^y\Psi(\vartheta,u)u^{-2}\dd u$,
con $\Psi=\sum_{n\le u}\Lambda(n)\vartheta(n)$.
Para el principal, la contribución de $u$ a esta integral es
$y\log y$. Sustituir las dos correcciones que definen $G_\vartheta$
y restar la integral total produce \eqref{eq:aux-integral-g}.
En ningún paso se introduce $L'/L(1,1)$.
\end{proof}

\subsection{Una regla de exponentes para el error relativo}

\begin{proposicion}[Separación entre conductor y argumento]
\label{prop:aux-delgada}
Supongamos \eqref{eq:aux-psi} con constantes uniformes y
\begin{equation}
 \lambda\le D x^\alpha(\log x)^b,\qquad
 \log y\ge c_0x^\beta,
 \quad 0<\alpha<\beta,\quad b\ge0,\quad D,c_0>0.
 \label{eq:aux-escala-general}
\end{equation}
Los parámetros de esta fórmula son fijos. Para cada
\begin{equation}
 0<r<\min(\beta-\alpha,\beta/2)
 \label{eq:aux-rango-r}
\end{equation}
hay un umbral de $x$, uniforme en el carácter y en $y$ admitidos,
tal que $|G_\vartheta(y)|/y\le e^{-x^r}$.
\end{proposicion}

\begin{proof}
La desigualdad \eqref{eq:aux-rango-r} permite elegir
\begin{equation}
 \max(\alpha/\beta,1/2)<\mu<1-r/\beta.
 \label{eq:aux-mu}
\end{equation}
Para $v\ge v_0=\log y$ y $x$ grande, las cotas de
\eqref{eq:aux-escala-general} implican $\lambda\le v^\mu$:
el factor fijo $D(\log x)^b$ se absorbe en la potencia
$x^{\beta\mu-\alpha}$, cuyo exponente es positivo.
También $\sqrt v\le v^\mu$ y $(v+\lambda)^2\le4v^2$.
Por \eqref{eq:aux-integral-g} y el cambio $v=\log u$,
\begin{equation}
 \frac{|G_\vartheta(y)|}{y}
 \le4C\int_{v_0}^\infty
                 v^2e^{-(c/2)v^{1-\mu}}\dd v.
 \label{eq:aux-cola-delgada}
\end{equation}
El cambio $w=v^{1-\mu}$ convierte la integral en la cola de
una potencia fija de $w$ multiplicada por $e^{-cw/2}$.
Para $w$ grande, esa potencia se acota por una constante por
$e^{cw/4}$, pues $\log w=o(w)$. Integrando el exponencial restante,
\[
 \frac{|G_\vartheta(y)|}{y}
 \le C_1e^{-c_1v_0^{1-\mu}}
 \le C_1e^{-c_2x^{\beta(1-\mu)}}.
\]
Como $\beta(1-\mu)>r$, se tiene
$\log C_1+x^r\le c_2x^{\beta(1-\mu)}$ a partir de un umbral fijo.
Esto da la conclusión con todas las uniformidades indicadas.
\end{proof}

En la aplicación tomamos $\alpha=\varepsilon/5$,
$\beta=\varepsilon/3$, $b=2$ y $r=\varepsilon/8$.
Para $y\ge\tfrac13e^{x^{\varepsilon/3}}$,
$\log y\ge\tfrac12x^{\varepsilon/3}$ a partir de un umbral.
El rango permitido es $r<2\varepsilon/15$, y
$\varepsilon/8<2\varepsilon/15$. La elección explícita
$\mu=61/100$ conserva las dos holguras
\begin{equation}
 \beta\mu-\alpha=\frac{\varepsilon}{300}>0,
 \qquad \beta(1-\mu)-r=\frac{\varepsilon}{200}>0.
 \label{eq:aux-holguras}
\end{equation}
La proposición produce $e^{-x^{\varepsilon/8}}$ después de aumentar
el umbral; en particular implica la cota con la reserva $1/2$ de
\eqref{eq:aux-meta-delgada}, utilizada por A93, Proposición~5.3.
El resultado se mantiene para cualquier $D$ fijo en
\eqref{eq:aux-conductor-delgado}. Una elección más amplia de $z'$
puede cambiar $D$; no se impone el valor $D=2$ a datos para los que
ese valor no haya sido comprobado.

\subsection{Compatibilidad con el cuerpo y con la ventana aritmética}

La condición sobre caracteres reales de la escala delgada debe
verificarse para los conductores del cuerpo que efectivamente se usa.
Si un carácter real primitivo tiene conductor divisor de $Q_0qd$,
se separan los factores $q,d$ y queda un carácter real primitivo de
conductor divisor de $8\operatorname{odd}(Q_0)$.
Cuando $\operatorname{odd}(Q_0)\mid M$, ese factor pertenece a la
familia de conductor divisor de $8M$ de \eqref{eq:aux-cajas}.
Si aparece $q$ o $d$, el torcido es no trivial y su caja excluye en
particular los ceros reales mayores que $399/400$.
La entrada formal \texttt{RealBox} conserva además las hipótesis sobre
el soporte primo del módulo del cuerpo; esas condiciones se comprueban
en la aplicación y no se sustituyen por una semejanza de notación.

Para combinar las ventanas, la Nota~\ref{sec:cohomologia} fija el
exponente $J$ de su selección antes de elegir la pareja. La
Proposición~\ref{prop:aux-pareja} admite ese $J$ fijo y produce
una pareja con estimaciones analíticas en todo $[Y,Y^J]$.
La entrada cohomológica se aplica a esa misma pareja admisible.
El reescalamiento de la Proposición~\ref{prop:par-ventana-comun}
conserva la pareja al pasar al intervalo común.

Quedan reconstruir las entradas uniformes recibidas en esta nota y
las realizaciones automorfas y de Galois que utilizan los mismos datos.
Los argumentos demostrados explican el segundo momento, la simultaneidad,
la selección dependiente y la absorción del error en escala delgada.
La Nota~\ref{sec:galois} estudia la entrada de Galois de A, §9,
pp.~42--46, y A93, §10, pp.~77--85: extracción, pegado y
compatibilidad local. El enlace integral con las congruencias
y las cotas cohomológicas conserva sus hipótesis adicionales.


\subsection{Concordancia localizada con Lean}
\label{subsec:aux-lean}

La octava lectura fuente conserva el espejo de
\texttt{CocoSgt/riemann-hypothesis-lean}, revisión
\nolinkurl{7cbf7a0c32fca1badf05931879a37c3f8a2b03b4}.
Comprende 34 archivos: 28 lecturas completas y seis selectivas,
con 3\,704 líneas leídas. Se verificaron los 34 hashes Git de los
bytes recuperados y se registró además SHA-256. Los rangos exactos,
las declaraciones y sus hipótesis están en la clave
\texttt{lean\_source\_review\_auxiliary\_primes} de
\texttt{REGISTRO.json}. No se ejecutó Lean ni se auditó la clausura
transitiva de imports.

La diferencia entre una proposición y su productor es decisiva:
\texttt{PairGood.lean}, líneas 28--30, define \texttt{PairEx}
como una proposición; \texttt{pairEx.lean}, líneas 112--169,
contiene el teorema \texttt{QH.pairEx : UpperTheory.Top.PairEx}
y su prueba. La Nota~\ref{sec:parametros} recibió esa proposición
como entrada. Esta lectura localiza su productor y desarrolla parte
de sus mecanismos, manteniendo identificadas las estimaciones que
siguen fuera del desarrollo.

\begin{center}
\small
\begin{tabularx}{\textwidth}{@{}p{.34\textwidth}X@{}}
\toprule
Archivo y declaración & Relación matemática comprobada \\
\midrule
\texttt{MeanSqAbs}, \texttt{MSqPairs} &
Retienen la clase excepcional y el término mixto
$2B\,RR$ junto a $2B^2$. \\
\texttt{PrimeSetTwoBounds}, \texttt{meanSqK} &
\texttt{MeanSq} admite toda secuencia acotada, lo que permite
cortar $\kappa$ por $n>t$ para la derivada. \\
\texttt{BadShare.bad\_share} &
Controla los testigos variables en $t$ mediante derivadas,
integración y bloques diádicos. \\
\texttt{CountBadCore}, \texttt{StagePow} &
Usan el conductor $fp$ para inyectar los primos descartados;
el ahorro relativo es $U^{-21/40}$. \\
\texttt{Stage1Succ}, \texttt{CharOf} &
Pasan a primitivas; la igualdad con el carácter original
se afirma en unidades. \\
\texttt{PairShare}, \texttt{PairExists} &
Fijan un umbral común y eligen sucesivamente $q$ y $d$
en la fibra correspondiente. \\
\bottomrule
\end{tabularx}
\end{center}

Hay dos diferencias de presentación que no deben ocultarse.
La función \texttt{sP} del código suma todas las potencias de un primo,
mientras que $A_\ell$ en \eqref{eq:aux-coeficientes} suma las impares
fuera del módulo. El código utiliza así un mayorante más grueso;
no se identifica esa suma con nuestra descomposición exacta.
Además, \texttt{BadShare} obtiene una proporción intermedia
\begin{equation}
 \frac{32C}{\tau}\frac{Y^{-\tau}}{1-2^{-\tau}}
 \label{eq:aux-diadas-lean}
\end{equation}
para las sumas que trata. El Lema~\ref{lem:aux-maximal}, siguiendo
A93, Lema~5.2, demuestra otra cota, con $Y^{-2\tau}$ y logaritmos,
para funciones absolutamente continuas bajo sus dos momentos.
Ambas pruebas dan la simultaneidad requerida; no se atribuye al
código literalmente la constante de la prueba escrita.

\begin{center}
\small
\begin{tabularx}{\textwidth}{@{}p{.34\textwidth}X@{}}
\toprule
Archivo y declaración & Paso y alcance \\
\midrule
\texttt{TwistPrim}, \texttt{BoxOfS} &
La identidad global para conductores coprimos permite
usar $G=-S-yT$ para la misma función. \\
\texttt{TailKExp.tailK} &
Conserva $Z=X^{400}$ y la holgura
$\varepsilon/(1600J)$ en la cola. \\
\texttt{Kappa4Loglog.kappa4} &
Pasa de $1000\log\log Q+C$ a $4\log Y$;
\texttt{pairEx} usa después $Y\le y$. \\
\texttt{PsiThin}, \texttt{Thin} &
Delimitan la fórmula explícita con el polo sustraído
y la cota en escala delgada. \\
\texttt{RealBox}, \texttt{thinH} &
Conservan las hipótesis de conductores y soporte primo
al trasladar la cota a los datos del cuerpo anfitrión. \\
\bottomrule
\end{tabularx}
\end{center}

Las seis lecturas selectivas son \texttt{PsiThin}, \texttt{twistK},
\texttt{MSqCount}, \texttt{MSqTwist}, \texttt{MSqMain} y
\texttt{MSqMainB}. La revisión no reconstruye todas las pruebas de
progresiones uniformes, Siegel, región sin ceros, densidad, contornos
ni fórmula explícita. Tampoco convierte la lectura de
\texttt{QH.pairEx} o \texttt{QH.thinH} en una nueva verificación
formal de sus dependencias. Este alcance precisa qué parte del puente
analítico se estudió y cuál permanece como entrada.

\subsection{Ejercicios resueltos}

\begin{ejercicio}\label{ej:aux-congruencias}
Tome $q=11$, $z=3$, $R_q=33$ y $d=263$.
Compruebe las congruencias de la segunda etapa y las condiciones
locales en $2$, $3$ y $q$. ¿Demuestra este ejemplo la existencia
en los intervalos exponenciales requeridos?
\end{ejercicio}

\paragraph{Solución del Ejercicio~\ref{ej:aux-congruencias}.}
$263$ es primo: no es divisible por $2,3,5,7,11,13$, los únicos
primos no mayores que su raíz cuadrada. Se tiene
\[
 8R_q-1=263,\qquad d\equiv7\pmod{16},\qquad
 d+1=264=8R_q.
\]
Por tanto $-d\equiv1\pmod3$ y $-d\equiv1\pmod{11}$;
$3$ y $11$ se descomponen en $\Q(\sqrt{-263})$.
Como $-263\equiv1\pmod8$, también se descompone $2$ en ese cuerpo.
En cambio $-11\equiv5\pmod8$, de modo que $2$ es inerte en
$\Q(\sqrt{-11})$. El ejemplo comprueba las congruencias y sus
consecuencias locales. No comprueba los intervalos exponenciales
ni las estimaciones analíticas, que requieren sus propios parámetros
y umbrales.

\begin{ejercicio}\label{ej:aux-proporcion}
Tres familias $F_{p,1},F_{p,2},F_{p,3}$ satisfacen los dos momentos
del Lema~\ref{lem:aux-maximal} con $C=1$ y
$\mathcal P\ne\varnothing$. Sean $h=1/2$, $\tau=1/4$ y
$Y\ge10^8$. Demuestre que al menos tres cuartas partes de los
índices cumplen las tres cotas en todo $[Y,X]$, cualquiera que sea
$X\ge Y$.
\end{ejercicio}

\paragraph{Solución del Ejercicio~\ref{ej:aux-proporcion}.}
La suma de las tres proporciones excepcionales es a lo sumo
\[
 12Y^{-1/2}\bigl(5\log(2Y)+8\bigr).
\]
La derivada de $Y^{-1/2}(5\log(2Y)+8)$ es
$Y^{-3/2}(1-\tfrac52\log(2Y))$, negativa para $Y\ge2$.
En $Y=10^8$, como $\log(2\cdot10^8)<20$, la cota es menor que
$12\cdot108/10^4=0.1296<1/4$. Los otros índices satisfacen las
tres desigualdades $|F_{p,j}(t)|\le\tfrac12t^{3/4}$ para todo
$t\in[Y,X]$. El extremo $X$ no entra en esta estimación.

\begin{ejercicio}\label{ej:aux-familia-creciente}
Sean $Y\ge3$, $X=2Y$ y
\[
 N=\left\lceil\frac{4\sqrt{2Y}}{\log(2Y)}\right\rceil,
 \qquad \mathcal P=\mathcal A=\{1,\ldots,N\},
 \qquad F_{p,\alpha}(t)=2t^{3/4}\mathbf1_{p=\alpha}.
\]
Compruebe los dos momentos con $C=1$. Muestre que cada condición
individual descarta una proporción que tiende a cero, pero ningún
índice cumple todas las condiciones con $h=1$, $\tau=1/4$.
\end{ejercicio}

\paragraph{Solución del Ejercicio~\ref{ej:aux-familia-creciente}.}
Para $\alpha$ fijo y $Y\le t\le2Y$,
\[
 \sum_p|F_{p,\alpha}(t)|^2=4t^{3/2}
 \le Nt\log(2t),
\]
pues $4\sqrt t/\log(2t)\le4\sqrt{2Y}/\log(2Y)\le N$.
Además $tF'_{p,\alpha}=\tfrac34F_{p,\alpha}$, de modo que
el segundo momento es $9/16$ del primero. Para cada $\alpha$
solo falla $p=\alpha$: su proporción es $1/N\to0$.
Pero cada $p$ falla para $\alpha=p$. Este ejemplo usa índices
abstractos, sin afirmar que sean primos. Demuestra que la finitud
a cada escala no reemplaza el control del número de condiciones
cuando la escala varía.

\begin{ejercicio}\label{ej:aux-fibras}
En el Lema~\ref{lem:aux-fibras}, tome
$\mathcal Q=\{0,1\}$, $\mathcal Q_0=\{0\}$,
una fibra $\mathcal D(0)$ de un único elemento, que es bueno, y una fibra
$\mathcal D(1)$ con $K\ge1$ elementos, ninguno bueno.
Compare la selección sucesiva con la medida uniforme sobre
el conjunto de todas las parejas.
\end{ejercicio}

\paragraph{Solución del Ejercicio~\ref{ej:aux-fibras}.}
Al elegir primero $q$ uniformemente y después un elemento de su
fibra uniformemente, la probabilidad de éxito es $1/2$:
el único elemento de la fibra de $0$ es bueno.
En la unión de todas las parejas hay $K+1$ elementos y solo uno
bueno, por lo que la proporción uniforme global es $1/(K+1)$.
Los tamaños distintos de las fibras explican la diferencia.
La existencia sucesiva necesita una fibra buena sobre un primer
índice retenido; no necesita independencia ni una medida uniforme
global sobre todas las parejas.

\begin{ejercicio}\label{ej:aux-holguras}
Use la Proposición~\ref{prop:aux-delgada} con
$\alpha=\varepsilon/5$, $\beta=\varepsilon/3$ y
$r=\varepsilon/8$. Compruebe el intervalo permitido para $r$
y las reservas que deja $\mu=61/100$.
\end{ejercicio}

\paragraph{Solución del Ejercicio~\ref{ej:aux-holguras}.}
Se tiene $\beta-\alpha=2\varepsilon/15$ y
$\beta/2=\varepsilon/6$, de modo que el extremo menor es
$2\varepsilon/15$. La diferencia con $r$ es
$2\varepsilon/15-\varepsilon/8=\varepsilon/120>0$.
Además $\alpha/\beta=3/5$ y $1-r/\beta=5/8$;
por tanto $3/5<61/100<5/8$. Finalmente
\[
 \beta\mu-\alpha=\frac{\varepsilon}{300},\qquad
 \beta(1-\mu)-r=\frac{\varepsilon}{200}.
\]
La primera reserva absorbe el factor fijo y los logaritmos del
conductor; la segunda absorbe las constantes y el prefactor de
la cola integral. Ambas son positivas para $\varepsilon$ fijo.
No dan un mismo umbral cuando $\varepsilon$ tiende a cero.
